\textbf{No photographed Indian bills.} The problem statement asks for generalisation to real Indian restaurant and canteen bills, and we did not have a labelled set of them. Everything we report about Indian-style layouts (charge lines, GST wording, rupee formats) is measured on bills our own generator draws, both the 100-bill test split and the 100 further bills used for the before and after comparison. Those scores show that the model has learned the output format and the charge vocabulary, and nothing more; they cannot tell how it reads a crumpled, folded or handwritten bill from a real mess. SROIE has no item labels, so the main task never saw that kind of image: on 8 real SROIE scans, with or without six image enhancements, it returned no items at all (it wrote the four SROIE fields instead, which the auxiliary task reads well: date 8 of 8, company 7, total 7, address 4). The app now uses that auxiliary reading to offer the shop name and total and asks for the items by hand. Whether real restaurant bills behave like these retail scans or like CORD's restaurant photos is exactly what real Indian bills would tell us.

\textbf{Evaluation caveats.} The metrics are our own implementation, not an official scorer. The EDA found many near-duplicate image pairs across the CORD and SROIE splits (likely the same shop's template), which we did not check by eye, so the held-out scores may be slightly generous. About 15\% of CORD labels do not satisfy their own arithmetic, which caps the arithmetic-consistency score and the CORD charge F1. The confidence analysis uses CORD and generated validation bills only; SROIE fields have no confidence. The checkpoint was selected on validation bills, so validation scores are a little optimistic compared with the test scores (CORD edit-distance score 0.958 against 0.930). The zero-shot comparison of larger models is not a comparison after training.

\textbf{Serving.} The reference app has no sign-in, rate limit or persistent storage, handles one bill per GPU at a time, and was load-tested on a single laptop GPU, not on a deployed server.

\subsection*{Future work}
\begin{itemize}
\item \textbf{Real bills.} Collect and label 60--150 photographed Indian bills, evaluate on at least 30 held-out ones, and fine-tune with the rest; this is the largest gap in the evidence.
\item \textbf{Running on the phone.} The fp16 model is 411 MB, too large to ship comfortably inside an app. Options are exporting the encoder and decoder to a mobile runtime (ONNX Runtime Mobile, Core ML, TensorFlow Lite), distilling into a smaller student, lowering the input resolution after measuring the accuracy cost, and quantising with a mobile toolchain. Our int8 and FP8 tests were run in eager PyTorch on a desktop GPU and do not predict how an NPU kernel behaves, so each of these needs its own accuracy check before it is trusted. A hybrid is also possible: a phone does capture and quality checks and the server reads the bill.
\item \textbf{Tilt and the retake prompt.} Extend the straightening estimator beyond 10 degrees and check it on real photos; retune the retake thresholds, which were set on CORD and SROIE photos and are over-cautious on our generated bills.
\item \textbf{Faster serving.} TensorRT or ONNX export of the encoder, a static key-value cache for the decoder, batching of concurrent requests, and a decision on whether to enable the compiled encoder after a larger test.
\item \textbf{Harder bills.} Hindi and other scripts, handwritten totals, bills longer than one photo (stitching several photos), partial-occlusion augmentation, and a fine-tuned vision-language model (for example a LoRA-trained Qwen3-VL-2B) as a challenger to compare accuracy against Donut.
\item \textbf{Product.} Saved groups, payment links, and an offline-capable installable web app.
\end{itemize}

\subsection*{Open-source compliance and reproducibility}
No closed or paid model or API is called anywhere in the extraction pipeline or the app; inference needs no network access. The base model, libraries, fonts and the CORD and SROIE data are open or publicly available, and \texttt{OPEN\_SOURCE.md} lists each component. Training code (\texttt{splitsnap/train.py} and the Kaggle runner), the split manifest builder, the evaluation scripts and the logs behind every figure in this report are in the repository.
